\documentclass{article}
\usepackage[utf8]{inputenc}
\usepackage[margin=0.75in]{geometry}
\usepackage[dvipsnames]{xcolor} % adding colour
\usepackage{amsmath} % adding math equations
\usepackage{natbib} % adding BibTeX support
\usepackage{hyperref}

\begin{document}
	\begin{center}

		\LARGE{\textbf{STA380 Project Proposal}} \\
        \vspace{1em}
        \Large{Permutation Test on ECG-Derived ``Age Gap'' Using the CODE-15\% Dataset (Zenodo 4916206)} \\
        \vspace{1em}
        \normalsize\textbf{Yanming Zhao, Meizhi Dong, Shijun Chen,\\Yaohe Wang, Hengchen Cao} \\
        \vspace{1em}
        \normalsize{University of Toronto, Mississauga}

	\end{center}

    \begin{normalsize}
    We propose a project under \textbf{Theme 3, Topic 11: Permutation Test}. We will use the publicly available \textbf{CODE-15\% ECG dataset} 
    hosted on Zenodo (record \textbf{4916206}) \citep{code15_zenodo_4916206}. The app will implement a two-sample permutation test comparing subjects with atrial fibrillation (AF) versus those without AF.

	\section{Project Topic}
    \textbf{Permutation Test (Advanced +2):} a two-sample permutation test for detecting differences between \textbf{AF} and \textbf{non-AF} groups on an ECG-derived continuous outcome.

	\section{Simulation vs. Dataset}
    This project is \textbf{dataset-based}. We will extract relevant variables from the dataset metadata file (e.g., \texttt{exams.csv}), including:
    \begin{itemize}
        \item \texttt{AF} (binary group label),
        \item \texttt{age} (chronological age),
        \item \texttt{nn\_predicted\_age} (neural-network predicted ``ECG age'').
    \end{itemize}
    Our main analysis does not require raw waveform files, but if time permits we will include an optional extension using simple waveform features as an additional outcome.

    \section{Simulation vs. Dataset}

    This project is \textbf{dataset-based}. We will use the publicly available
    \textbf{CODE-15\% ECG dataset} hosted on Zenodo (record 4916206) 
    \citep{code15_zenodo_4916206}.

    \subsection*{Dataset Description}

    The CODE-15\% dataset is a large-scale clinical electrocardiogram (ECG) database
    consisting of approximately 345,779 ECG exams from 233,770 patients collected
    between 2010 and 2016. Each exam contains demographic information (e.g.,
    chronological age and sex), automated diagnostic labels (e.g., atrial
    fibrillation (AF), right bundle branch block, and other ECG abnormalities),
    and a neural-network–predicted ECG age derived from a deep learning model.

    Although the raw ECG waveforms (12 leads, 4096 time points at 400 Hz sampling rate)
    are available, our primary analysis focuses on structured metadata in
    \texttt{exams.csv}. In particular, we will use:

    \begin{itemize}
        \item \texttt{AF}: binary indicator for atrial fibrillation,
        \item \texttt{age}: chronological age,
        \item \texttt{nn\_predicted\_age}: deep neural network predicted ECG age.
    \end{itemize}

    \subsection*{Outcome Definition}

    We define the continuous outcome variable:

    \[
    g = \texttt{nn\_predicted\_age} - \texttt{age},
    \]

    which we refer to as the \textit{Age Gap}. A positive value ($g>0$) indicates
    that the model predicts an ECG-derived physiological age greater than the
    chronological age, while a negative value ($g<0$) indicates the opposite.
    This provides a clinically meaningful continuous variable suitable for
    two-sample comparison.

    \subsection*{Statistical Considerations}

    Because multiple ECG exams may correspond to the same patient
    (\texttt{patient\_id}), observations may not be strictly independent.
    To address this issue, our default implementation will restrict the analysis
    to one exam per patient (e.g., earliest or randomly selected exam).
    As an optional extension, we will implement a patient-level block permutation
    procedure to preserve within-patient dependence structure.

	\section{Project Details}
    \textbf{Research question.} Do patients with \textbf{AF} and \textbf{non-AF} have statistically significant differences in an ECG-derived ``age gap''?

    \vspace{0.5em}
    \noindent \textbf{Groups (two samples).}
    \begin{itemize}
        \item Group 1: \texttt{AF == TRUE} (AF)
        \item Group 2: \texttt{AF == FALSE} (non-AF)
    \end{itemize}

    \noindent \textbf{Outcome variable.} We define the continuous outcome (``Age Gap''):
    \[
        g = \texttt{nn\_predicted\_age} - \texttt{age}.
    \]
    Intuitively, $g>0$ means the ECG-based model predicts an ``older'' physiological age than chronological age, while $g<0$ suggests the opposite.

    \vspace{0.5em}
    \noindent \textbf{Hypotheses.}
    \begin{itemize}
        \item $H_{0}$: The distribution of $g$ is the same for AF and non-AF (group label is exchangeable).
        \item $H_{1}$: The distributions differ (mean/median/overall shape differs).
    \end{itemize}

    \noindent \textbf{Test statistic (user-selectable).} To demonstrate flexibility of permutation testing, users can choose:
    \begin{enumerate}
        \item Mean difference: $T = \bar g_{\text{AF}} - \bar g_{\text{nonAF}}$
        \item Median difference: $T = \mathrm{med}(g_{\text{AF}}) - \mathrm{med}(g_{\text{nonAF}})$
        \item Kolmogorov--Smirnov statistic: $T = \sup_{x} | \hat F_{\text{AF}}(x) - \hat F_{\text{nonAF}}(x) |$
    \end{enumerate}

    \noindent \textbf{Permutation procedure.} Let $T_{\text{obs}}$ be the observed statistic computed from the original labels.
    \begin{enumerate}
        \item Fix the observed $g$ values and randomly permute AF labels across observations.
        \item For $b=1,\dots,B$, compute the permuted statistic $T_b$.
        \item Compute a two-sided p-value:
        \[
            p = \frac{1 + \sum_{b=1}^{B} \mathbf{1}\left(|T_b| \ge |T_{\text{obs}}|\right)}{B+1}.
        \]
        \item At level $\alpha$, report \textbf{reject} or \textbf{fail to reject} $H_0$.
    \end{enumerate}

    \noindent \textbf{Independence / repeated exams.} Since the dataset may contain multiple exams per patient, we will explicitly address non-independence by providing:
    \begin{itemize}
        \item \textbf{Default (recommended):} restrict to \textbf{one exam per patient} (user may select ``earliest'' or ``random'' exam per \texttt{patient\_id}).
        \item \textbf{Optional extension:} a \textbf{block permutation} approach at the patient level (permute group labels at the \texttt{patient\_id} level), and compare results to the default approach.
    \end{itemize}

    \noindent \textbf{Outputs provided.}
    \begin{itemize}
        \item A histogram/density plot of $\{T_b\}_{b=1}^B$ with a vertical line at $T_{\text{obs}}$
        \item The permutation p-value and a clear reject/fail-to-reject statement
        \item A summary table including group sizes and descriptive statistics of $g$ (mean, median, SD)
        \item (Optional) sensitivity comparison across statistics and sampling modes (per-patient vs exam-level)
    \end{itemize}

	\section{User Inputs (Shiny Components)}
    Users will be able to modify the following:
    \begin{enumerate}
        \item Random seed
        \item Number of permutations $B$ (e.g., 500 / 1000 / 5000)
        \item Significance level $\alpha$ and one-sided vs two-sided testing
        \item Choice of test statistic (mean difference / median difference / KS)
        \item Data handling option for repeated exams:
            \begin{itemize}
                \item one exam per patient (earliest vs random),
                \item (optional) patient-level block permutation.
            \end{itemize}
        \item Display options (toggle):
            \begin{itemize}
                \item permutation distribution plot,
                \item group summary table,
                \item brief assumptions/caveats panel.
            \end{itemize}
    \end{enumerate}

\end{normalsize}

\nocite{*}
\bibliographystyle{abbrvnat}
\bibliography{bibliography}

\end{document}
